\begin{result}{E.2.11}{The prescribed CM comparison has real and rational abelian lifts}{Proved auxiliary result; full integral gate retained}
Let $F_{\rm CM}\subset\Qb$ be the compositum of $K^{\rm ab}$ for all
finite CM fields $K$, and take $\sigma\in\Gal(\Qb/F_{\rm CM})$.
Let $\mathrm{AV}(\Q)$ be the category of pure abelian motives with
absolute Hodge correspondences, including Artin motives and Tate
objects. Write
\[
 \pi:G_{\mathrm{AV},B}\longrightarrow G_{\mathrm{CM},B},
 \qquad D=\ker\pi
\]
for restriction of its rational Betti group to the CM/Artin subcategory.
The prescribed CM comparison $g_\sigma$ of E.2.6 defines a point of
$G_{\mathrm{CM},B}(\Q)$. Its fibre
\[
 X_\sigma=\pi^{-1}(g_\sigma)
 \tag{E.2.11a}
\]
has both a real point and a rational point.

The real point can be chosen to preserve Hodge structures and labelled
polarizations. The rational point is a tensor comparison in the
absolute-Hodge abelian category. This assertion supplies neither
integrality of that rational point nor its equality with the chosen
real Hodge point. The full pure integral question (E.2.3d) remains open
at this step.
\end{result}
\begin{proof}
\emph{The prescribed CM point and its morphisms.}
Theorem~3.15 of
\cite[p.~41 in the corrected author version]{GenMilneShimura1994}
gives, for a CM motive $M$ and a representative $\widetilde r(\sigma)$
of its reciprocity class, a unique isomorphism
$\alpha_M:M\to{}^\sigma M$ carrying every labelled Hodge tensor $s$
to $\sigma(s)$, with
\[
 \sigma x=(\alpha_M)_f\bigl(\widetilde r(\sigma)x\bigr)
                  \qquad(x\in\omega_f(M)).
 \tag{E.2.11b}
\]
Every relevant reflex field $E$ is CM or totally real. In the totally
real case $E(i)$ is CM and $E^{\rm ab}E(i)/E(i)$ is abelian, so
$E^{\rm ab}\subset E(i)^{\rm ab}\subset F_{\rm CM}$.
The containment is immediate when $E$ is CM. Thus the reciprocity
class is trivial, and the representative $1$ in (E.2.11b) gives the
precise maps of E.2.6. This verifies their naturality for all
absolute Hodge correspondences as well as algebraic ones.
Their Artin components are $e_x\mapsto e_{\sigma x}$ and their
component on the one Tate line is the identity.

\emph{A finite comparison with the prescribed CM restriction.}
Consider finitely many abelian varieties over $\Qb$, finitely many
labelled absolute Hodge tensor maps, chosen polarizations, and a
finite set of CM objects on which the comparison is prescribed.
Take one product $A$ of the abelian generators, including a product
$B$ of CM abelian generators for those CM objects. Include the factor
projectors, the polarization tensors, and the one Tate object.
Let $V=H_1(A(\C),\Q)$, $\Lambda=H_1(A(\C),\Z)$, and
$G=\operatorname{MT}(V,h)$. Enlarge the indexed Hodge tensors to
cut out $G$, retaining all the original labels.

The $G(\R)$-conjugacy class $X$ of $h$ gives its Hodge-type Shimura
datum. The exact moduli statements used here are
\cite[Example~3.1, p.~36; Lemma~3.29 and Theorem~3.31,
pp.~45--46]{GenMilneShimura1994}: the canonical model parametrizes
the indexed absolute Hodge tensors and finite-adic level structure,
and conjugating over the reflex field conjugates those labelled
data. The integral local system is explicit in the same source:
conditions (2.7.1)--(2.7.3), Proposition~2.8, Example~2.11 and
Theorem~2.12, pp.~23--25, construct the constant lattice on $X$
and its quotient on $\Gamma(N)\backslash X$. Take $N$ sufficiently
divisible, with $3\mid N$. Example~3.7, p.~38, identifies this with
the $K(N)$-level description. The level orbit determines
$\eta(\Lambda\otimes\widehat\Z)$; intersection with the rational
realization recovers the actual integral lattice. Thus no unspecified
lattice on an isogeny-class representative is substituted.

The reflex field $E(G,X)$ is a subfield of a CM field by the proof
of \cite[Theorem~3.33, p.~46]{GenMilneShimura1994}. Such a subfield
is CM or totally real: in a CM Galois closure complex conjugation is
central, hence preserves each intermediate field, acting either
trivially or as its CM involution. The preceding calculation gives
$E(G,X)^{\rm ab}\subset F_{\rm CM}$.
Deligne's component reciprocity theorem
\cite[Section~2.6.1 and Theorem~2.6.3, p.~46 of the translation]{GenDeligneShimura1979}
therefore puts the point defined by $A$ and its $\sigma$-conjugate
in the same connected component at this level.
Parallel transport along a path between them supplies an integral
map
\[
 T_0:H_1(A(\C),\Z)\xrightarrow{\sim}
                     H_1({}^\sigma A(\C),\Z)
 \tag{E.2.11c}
\]
preserving every labelled tensor and compatible with $\sigma$ on
level $N$. Its Tate action is the identity: on the integral rank-one
Tate local system it is a unit congruent to $1$ modulo $3$.

The CM restriction of (E.2.11c) is exactly $\alpha_B$.
Indeed the representation of $G$ on $H_1(B)$ has torus image;
the corresponding Hodge homomorphism is constant under conjugation
along the connected domain $X^+$. Hence $T_{0,B}$ is an integral
Hodge isometry. Both $T_{0,B}$ and $\alpha_B$ induce $\sigma$ on
$B[N]$. Their quotient is a polarized integral Hodge automorphism
that is the identity modulo $3$, and such an automorphism is the
identity. To check the last assertion, it preserves the positive
definite form obtained from the polarization and Weil operator,
so its integral matrix lies in a finite group. But
$1+3M_r(\Z_3)$ is torsion-free. A matrix $U=I+3^sV$, $s\ge1$,
with $V\not\equiv0\pmod3$ cannot have prime order $p\ne3$ by
binomial expansion modulo $3^{s+1}$; for $p=3$,
\[
 U^3=I+3^{s+1}\bigl(V+3^sV^2+3^{2s-1}V^3\bigr)
\]
has the parenthesis nonzero modulo $3$. Taking a prime-order power
excludes every nontrivial finite-order element.

The map $T_0$ need not preserve the Hodge structure of $A$.
After trivializing along the path, $h$ and
$T_0^{-1}h_{{}^\sigma A}T_0$ belong to the same $X^+$.
The map
$G^{\rm der}(\R)^+\to G^{\rm ad}(\R)^+$
is surjective: its differential is an isomorphism, so its image is
open in the connected target. Choose $d$ in the derived group
carrying the first Hodge homomorphism to the second. Then
\[
 T=T_0d
 \tag{E.2.11d}
\]
is a real Hodge isometry preserving all the indexed tensors.
The derived group acts trivially on CM representations and on
Tate, so $T_B=\alpha_B$. It preserves polarizations because their
similitude characters factor through the abelian quotient.
This proves finite real solvability with the prescribed CM map.

\emph{Coherence and Artin blocks.}
For every abelian variety $A/\Qb$, choose a polarization and let
$C_A$ be the set of real Hodge isometries from $H_1(A)$ to
$H_1({}^\sigma A)$ for that polarization and its conjugate.
Each $C_A$ is nonempty by the preceding construction and compact:
after choosing one isometry it is a torsor under a closed subgroup
of the orthogonal group of the associated positive definite form.
In $\prod_A C_A$, impose all naturality, tensor, unit, duality,
projector and labelled-polarization equations, together with
$T_B=\alpha_B$ for every CM generator. These are closed conditions,
each involving finitely many coordinates. Every finite collection
has a solution by (E.2.11d), applied to its finite product of
generators. Compactness and the finite intersection property give
a simultaneous solution.

Extend it by tensor operations and exterior powers; restrict to
the image of each projector. Naturality for the transition
isomorphisms between presentations makes the resulting comparison
well-defined on every abelian motive over $\Qb$.
For descent to $\Q$, send $e_u$ to $e_{\sigma u}$ on each finite
Artin embedding set. On the underlying $\Q$-scheme of an object
over a finite field $L$, send its $u$-summand to its $\sigma u$-summand
by the constructed map. The Artin module action selects these
blocks, and $\sum_u e_u\otimes e_u$ selects the relative tensor
product, exactly as in E.2.3. Each correspondence equation holds
on every geometric block and hence on their direct sum. Restriction
of embedding sets gives compatibility under changing $L$.
On $\Q$-defined objects, the canonical identification with their
$\sigma$-conjugates yields
$g_{\R}\in G_{\mathrm{AV},B}(\R)$ with
$\pi(g_{\R})=g_\sigma$. This proves $X_\sigma(\R)\ne\varnothing$.
No geometric real-Frobenius intertwining has been imposed;
E.2.7 gives its separate obstruction for the quartic places.

\emph{The rational point.}
For every finite prime $\ell$, Galois transport by $\sigma$ on
the etale realization, carried through Betti--etale comparison,
gives $g_\ell\in G_{\mathrm{AV},B}(\Q_\ell)$.
Naturality includes absolute Hodge correspondences through their
etale realizations. Equation (E.2.11b) with representative $1$
proves $\pi(g_\ell)=g_\sigma\otimes\Q_\ell$.
Thus the actual fibre (E.2.11a) has a point over every completion.

The proof of \cite[Proposition~3.12 and Theorem~3.13,
p.~26 of the author version]{GenMilnePeriods2004}
identifies $D$ as a countable product of semisimple algebraic
factors whose geometric simple factors have covering index $1$
or $2$. Proposition~3.5, pp.~23--24, proves the all-local Hasse
injection for those factors, and Corollary~1.14, p.~9, identifies
$H^1$ of their countable product with the product of the $H^1$
sets. Consequently
\[
 H^1(\Q,D)\longrightarrow\prod_v H^1(\Q_v,D)
\]
is injective. The class of $X_\sigma$ is zero in every target
factor by the points just constructed. It is therefore zero
over $\Q$, giving $X_\sigma(\Q)\ne\varnothing$.
This uses the proved product structure; it does not infer a
compatible lift from unrelated points of an inverse system.

The rational point obtained at this last step is not required to
equal $g_{\R}$ after scalar extension. In particular the construction
has proved a rational abelian comparison with the prescribed CM
restriction, while the simultaneous lattice conditions of E.2.9
and the full pure-motive requirement in (E.2.3d) still require their
own proofs.
\end{proof}
